\begin{textsample}{NEG Dim 3 – global\_south\_2025\_06 – Score -1.00 – global\_south\_2025\_06/125934659.txt}  \label{ex:f3_neg_270}
Sharad Pawar says India ' s stance on Gaza not in line with traditional position NCP founder Sharad Pawar on Monday questioned the central government ' s"neutral"stance about events in the Gaza strip, and said India ' s present stance of ' neutrality ' is not in line with its traditional position.
Statesman News Service Mumbai June 16, 2025 8:02 pm Photo : IANS NCP founder Sharad Pawar on Monday questioned the central government ' s"neutral"stance about events in the Gaza strip, and said India ' s present stance of ' neutrality ' is not in line with its traditional position.
He made the statement while addressing NCP workers in Mumbai ' s Ghatkopar suburb about the soon to be announced local civic body elections in Maharashtra.
Advertisement The Sharad Pawar-led NCP is part of the Maha Vikas Aghadi ( MVA ) alliance with Congress and the Uddhav Thackeray-led Shiv Sena.
Advertisement Speaking to NCP workers in Ghatkopar, Sharad Pawar referred to India abstaining from voting on a United Nations @ @ @ @ @ @ @ @ @ @ had called for an"immediate, unconditional and permanent"ceasefire in Gaza Strip, due to the humanitarian crisis in the region.
Pawar said that India had abstained from voting on the United Nations General Assembly resolution introduced by Spain and backed by 149 nations, which was overwhelmingly adopted by the 193-member UN General Assembly.
While 12 countries opposed the resolution, 19 nations including India, chose to abstain.
Pawar said that"Israel is doing good work in the field of agriculture, but on the other hand, what is happening in Gaza, where innocent children and people have lost their lives, their homes have been destroyed.
While all this is happening, many countries in the world chose to remain silent"."Whether it was during late Prime \textbf{Minister} Indira Gandhi ' s time, or Jawaharlal Nehru ' s time, or later, India always took a different stand on this issue.
India ' s stand always protected humanity.
That was India ' s image.
But just yesterday, in the @ @ @ @ @ @ @ @ @ @, and instead of taking a stand against aggression, India chose neutrality.
This is not India ' s policy."Unfortunately, those who have the reins of the country today do not present a clear and clean policy on such matters before the world.
Because of this, a large section of the world is developing, and will develop, a misunderstanding about India.
This is the situation we are seeing today,"Pawar said.
Speaking about the recent Pahalgam terror attack in Kashmir Pawar said that the central government should have taken a firm stand."Government ' s action should have been against the perpetrators who came with guns and killed innocent people, but that did not happen".
Pawar went on to raise concerns about the credibility of the election process itself, saying that people have doubts about the election process and even Rahul Gandhi has written about it."There is doubt in people ' s minds regarding the electoral machinery,"Pawar said.
He said that local NCP @ @ @ @ @ @ @ @ @ @ body polls independently or in alliance with other parties.
He assured the NCP workers that the party would fully support whatever decision they stood after studying the local political situation."In areas going to polls, local NCP workers and leaders should sit together and decide what ' s best.
If you choose to contest independently, the party is with you.
If you feel there should be discussion with other parties, the party will support that as well,"Pawar told NCP workers.

% matched lemmas: minister
\end{textsample}
